% This file is generated from the corresponding Obsidian note.
% Source: Teoricos/GS - Teo11.md
\chapter{Factorización y subvariedades incrustadas}
\label{ch:teo-11}
\section{Subvariedades: factorización e incrustadas}
\begin{bookproposition}{Existencia de cartas dado diferenciales}
Sea $M$ variedad suave de dimensión $m$ y sea $k\ge m$. Si

\[
\{y_1,\ldots,y_k\}
\]
es un conjunto de funciones suaves definidas en un abierto $U$ de $M$ en $p$ tal que

\[
\{(dy_1)_p,\ldots,(dy_k)_p\}
\]
generan $T_p^\ast M$, entonces existe un subconjunto

\[
\{i_1,\ldots,i_m\}\subseteq\{1,\ldots,k\}
\]
tal que

\[
\psi=(y_{i_1},\ldots,y_{i_m})
\]
define una carta suave en algún abierto de $p$.

\end{bookproposition}
\begin{bookproof}{}
\begin{enumerate}
\item Como $\{(dy_j)_p:j\in\{1,\ldots,k\}\}$ genera $T_p^\ast M$, entonces contiene una base
\end{enumerate}
\[
\{(dy_{i_1})_p,\ldots,(dy_{i_m})_p\}
\]
de $T_p^\ast M$. 2. Por \hyperref[blk:2b23b2]{Teórico 10}, la aplicación

\[
\psi=(y_{i_1},\ldots,y_{i_m})
\]
define una carta suave en cierto abierto de $p$.

\bookqed
\end{bookproof}
\phantomsection\label{blk:existencia-cartas-dado-diferenciales}
\phantomsection\label{blk:49431b}
\begin{bookcorollary}{}
Sean $M$ y $N$ variedades suaves de dimensión $m$ y $n$, respectivamente. Sea

\[
F:M\to N
\]
suave y supongamos que

\[
(dF)_p:T_pM\to T_{F(p)}N
\]
es inyectiva, y por tanto $m\le n$. (Osea $F$ es inmersion) Si $(V,\varphi=(x_1,\ldots,x_n))$ es una carta suave de $N$ en $F(p)$, entonces existen

\[
i_1,\ldots,i_m\in\{1,\ldots,n\}
\]
tales que

\[
\psi=(x_{i_1}\circ F,\ldots,x_{i_m}\circ F)
\]
define una carta suave de $M$ en cierto abierto de $p$.

\end{bookcorollary}
\begin{bookproof}{}
\begin{enumerate}
\item Para usar el \hyperref[blk:existencia-cartas-dado-diferenciales]{resultado anterior}, alcanza con ver que
\end{enumerate}
\[
\{d(x_1\circ F)_p,\ldots,d(x_n\circ F)_p\}
\]
genera $T_p^\ast M$. 2. Como

\[
(dF)_p:T_pM\to T_{F(p)}N
\]
es inyectiva, entonces

\[
(dF)^\ast_{F(p)}:T_{F(p)}^\ast N\to T_p^\ast M
\]
es sobreyectiva por \hyperref[blk:dual-inyectiva-sobreyectiva]{dual de una transformacion inyectiva}. 3. Como $(V,\varphi=(x_1,\ldots,x_n))$ es una carta suave de $N$ en $F(p)$ sucede que

\[
\left\{(dx_1)_{F(p)},\ldots,(dx_n)_{F(p)}\right\}
\]
es base de $T_{F(p)}^\ast N$ por \hyperref[blk:base-dual-coordenadas]{base dual de las coordenadas}. Entonces sus imágenes por $(dF)^\ast_{F(p)}$ generan $T_p^\ast M$ por el paso 2. 4. Pero

\[
(dF)^\ast_{F(p)}\big((dx_i)_{F(p)}\big)=d(x_i\circ F)_p.
\]
Luego los $d(x_i\circ F)_p$ generan $T_p^\ast M$, como queríamos.

\bookqed
\end{bookproof}
\phantomsection\label{blk:c4d494}
\section{Lema de factorización MUY IMPORTANTE}
\begin{bookremark}{Sobre restriccion codominio}
Tenemos la siguiente situación:

\begin{itemize}
\item $M,N$ y $P$ son variedades suaves.
\item $\iota:P\to N$ es una subvariedad de $N$, no necesariamente incrustada.
\item $F:M\to N$ es una función suave con
\end{itemize}
\[
F(M)\subseteq P.
\]
Se quiere restringir el codominio de $F$ a $P$, es decir, preguntar cuándo existe una función suave

\[
\widehat F:M\to P
\]
tal que

\[
F=i\circ\widehat F.
\]
En general no es cierto que restringir el codominio nos deje automáticamente una función suave.

\end{bookremark}
\begin{bookexample}{Ejemplo donde co-restringir no da suave}
Consideremos el ocho. Tenemos

\[
f:\mathbb R\to\mathbb R^2,\qquad t\mapsto \sin t(\cos t,1),
\]

\bookimage{assets/pasted-image-20260429114226.png}{Pasted image 20260429114226}
 Tomemos $M=(\text{Ocho},f^{-1}|_{(0,2\pi)})$ y $N=\mathbb{R}^2.$ También tomamos $P=(\text{Ocho},f^{-1}|_{(-\pi,\pi)})$ y $F=i:M\to N.$ Claramente la imagen de $F$ es el conjunto dado, en particular cae en $P$.\\
La restricción del codominio de $F$ a $P$,

\[
\widehat{F}:(\text{Ocho},f^{-1}|_{(0,2\pi)})\to (\text{Ocho},f^{-1}|_{(-\pi,\pi)})
\]
es la función identidad, $\widehat F=\operatorname{Id}$. (Que $\widehat{F}=\operatorname{Id}$ sale similar a \hyperref[blk:8ed2bd]{Teórico 9} creo, chequear)\\
Pero la identidad no puede ser continua, por que las topologias son distintas y no compatibles o usando sucesiones (Similarmente podemos afirmar que son subvariedades de $\mathbb{R}^{2}$ no equivalentes)

\end{bookexample}
\begin{bookproposition}{Lema de factorizacion}
Sea $F:M\rightarrow N$ suave. Y sea $(P,g)$ subvariedad de $N$ tal que $F(M)\subseteq g(P)$. Ademas sea $\widehat{F}:M\rightarrow P$ la funcion $\widehat{F}(x)=y$ si $F(x)=g(y)$ la cual es la unica funcion de $M$ en $P$ que hace conmutar el diagrama 
\bookimage{assets/pasted-image-20260429160319.png}{Pasted image 20260429160319}
 Entonces si $\widehat{F}$ es continua es suave

\end{bookproposition}
\begin{bookproof}{}
\begin{enumerate}
\item Sea $p\in M$ y veamos que $\widehat F$ es suave en un entorno de $p$.
\item Sea $q=\widehat F(p)$, es decir, $q$ es el único elemento de $P$ tal que
\end{enumerate}
\[
F(p)=g(q).
\]
\begin{enumerate}[start=3]
\item Sea $(V,\psi=(y_1,\ldots,y_n))$ una carta suave de $N$ alrededor de $F(p)$.
\item Como $g:P\to N$ es subvariedad, tenemos
\end{enumerate}
\[
(dg)_q:T_qP\to T_{g(q)}N
\]
es inyectiva, y por \hyperref[blk:c4d494]{Corolario}, existe un entorno $W$ de $q$ e índices $i_1,\ldots,i_k\in\{1,\ldots,n\}$, con $k=\dim P$, tal que

\[
(W,\varphi=(y_{i_1}\circ g,\ldots,y_{i_k}\circ g))
\]
es carta suave de $P$. 5. Sea

\[
\pi:\mathbb R^n\to\mathbb R^k,\qquad (s_1,\ldots,s_n)\mapsto (s_{i_1},\ldots,s_{i_k}),
\]
la cual es suave. 6. La ventaja de introducir esta función es que podemos expresar la carta $(W,\varphi)$ como

\[
(W,\pi\circ\psi\circ g).
\]
\begin{enumerate}[start=7]
\item Como $\widehat F$ es continua, tenemos
\end{enumerate}
\[
U=\widehat F^{-1}(W)
\]
es un abierto de $M$ que contiene a $p$. 8. Para ver que $\widehat{F}$ es suave en $p$ alcanza ver que $\widehat{F}|_{U}$ es suave y para esto componemos con la carta $\varphi$ (basta componer solo a un lado por que si vemos eso luego dada carta $\psi$ tendremos que $\varphi\circ\widehat{F}|_{U}\circ\psi$ es composicion de suaves)

\[
\varphi\circ\widehat{F}|_{U}=\pi\circ\psi\circ g\circ\widehat{F}|_{U}=\pi\circ\psi\circ F|_{U}
\]
esto ultimo vale por que $g(\widehat{F}|_{U}(x))=g(y)=F(x)$ por definicion 9. $\pi$ es suave por definicion $\psi\circ F|_{U}$ es suave por que $F$ es suave entonces $\varphi\circ\widehat{F}|_{U}$ resulta suave. Luego $\widehat{F}$ es suave

\bookqed
\end{bookproof}
\phantomsection\label{blk:1258d0}
\begin{bookremark}{}
Tenemos visto que restringir el codominio es peligroso para justificar que una función es suave.

En la práctica, la proposición anterior es una gran herramienta para justificar suavidad de una función

\[
\widehat F:M\to P
\]
si $P$ está metida como una variedad ambiente $N$ y la función

\[
F:M\to N
\]
que se obtiene al componer con la inclusión es más fácil de estudiar.

\end{bookremark}
\begin{bookproposition}{Corestringir a un subespacio con topología heredada preserva continuidad}
Sea $F:M\to N$ continua y sea $A\subseteq N$ con la topología heredada de $N$. Si $F(M)\subseteq A$, entonces la corestricción

\[
F_A:M\to A,\qquad F_A(x)=F(x)
\]
es continua.

\end{bookproposition}
\begin{bookproof}{}
10. Sea $W\subseteq A$ abierto. 11. Como $A$ tiene la topología heredada de $N$, existe un abierto $V\subseteq N$ tal que

\[
W=A\cap V.
\]
\begin{enumerate}[start=12]
\item Queremos probar que $F_A^{-1}(W)$ es abierto en $M$.
\item Tenemos
\end{enumerate}
\[
F_A^{-1}(W)=F_A^{-1}(A\cap V).
\]
\begin{enumerate}[start=14]
\item Como $F_A(M)\subseteq A$, para todo $x\in M$ se tiene
\end{enumerate}
\[
F_A(x)\in A.
\]
\begin{enumerate}[start=15]
\item Entonces
\end{enumerate}
\[
F_A(x)\in A\cap V \iff F_A(x)\in V.
\]
\begin{enumerate}[start=16]
\item Como $F_A(x)=F(x)$, se sigue que
\end{enumerate}
\[
F_A^{-1}(A\cap V)=F^{-1}(V).
\]
\begin{enumerate}[start=17]
\item Por lo tanto
\end{enumerate}
\[
F_A^{-1}(W)=F^{-1}(V).
\]
\begin{enumerate}[start=18]
\item Como $F:M\to N$ es continua y $V\subseteq N$ es abierto, entonces $F^{-1}(V)$ es abierto en $M$.
\item Luego $F_A^{-1}(W)$ es abierto en $M$.
\item Como esto vale para todo abierto $W\subseteq A$, concluimos que $F_A:M\to A$ es continua.
\end{enumerate}
\bookqed
\end{bookproof}
\begin{bookcorollary}{Lema de factorización para subvariedades incrustadas}
Sean $M,N,P$ variedades suaves. Sea $(P,g)$ subvariedad \textbf{incrustada} de $N$, es decir, $(P,g)$ es subvariedad y $g$ es un homeomorfismo entre $P$ y $g(P)$ con la topología heredada de $N$. Sea $F:M\to N$ suave tal que $F(M)\subseteq g(P)$. Entonces la única función $\widehat F:M\to P$ que hace conmutar el diagrama 
\bookimage{assets/pasted-image-20260429171650.png}{Pasted image 20260429171650}
 es continua (por lo tanto) y es suave.

\end{bookcorollary}
\begin{bookproof}{}
\begin{enumerate}
\item Como $F(M)\subseteq g(P)$, para cada $x\in M$ existe un único $y\in P$ tal que
\end{enumerate}
\[
F(x)=g(y).
\]
\begin{enumerate}[start=2]
\item Definimos
\end{enumerate}
\[
\widehat F:M\to P,\qquad \widehat F(x)=y \iff F(x)=g(y).
\]
\begin{enumerate}[start=3]
\item Entonces, por definición,
\end{enumerate}
\[
g\circ \widehat F=F.
\]
\begin{enumerate}[start=4]
\item Como $(P,g)$ es una subvariedad incrustada, la aplicación
\end{enumerate}
\[
g:P\to g(P)
\]
es un homeomorfismo, donde $g(P)$ tiene la topología heredada de $N$. 5. Como $F:M\to N$ es suave, en particular es continua. 6. Como $F(M)\subseteq g(P)$, podemos mirar $F$ como una función

\[
F:M\to g(P).
\]
pero ahora esta función sigue siendo continua porque $g(P)$ tiene la topología heredada de $N$. 7. Luego

\[
\widehat F=g^{-1}\circ F,
\]
con $g^{-1}:g(P)\to P$ es continua (Por ser $g$ homeo) . 8. Por lo tanto $\widehat F$ es continua. 9. Por el lema de factorización, si $\widehat F$ es continua entonces $\widehat F$ es suave. 10. Luego

\[
\widehat F:M\to P
\]
es suave.

\bookqed
\end{bookproof}
\phantomsection\label{blk:300685}
\begin{bookremark}{}
Supongamos que $(A,i)$ es subvariedad de $M$, con $i$ la inclusión, \textbf{ojo: no necesariamente incrustada}. Como $i$ es suave, en particular continua, entonces si $U$ es abierto de $M$, tenemos

\[
i^{-1}(U)=U\cap A
\]
es abierto de $A$. Por tanto, la topología que trae $A$, por ser variedad, es posiblemente más fina que su topología heredada del ambiente $M$. Un ejemplo es el ocho $\left(\text{ocho},f|_{(0,2\pi)}^{-1}\right)$, con $f:\mathbb R\to\mathbb R^2$,

\[
f(t)=\sin(t)(\cos(t),1).
\]
El conjunto

\[
I=f\left(\frac{\pi}{2},\frac{3\pi}{2}\right)
\]
es un abierto de esta variedad por la topologia que le damos que es la inducida por el homeo $f|_{(0,2\pi)}$, pero no puede obtenerse como un abierto relativo del ocho respecto a la topología de $\mathbb R^2$. 
\bookimage{assets/pasted-image-20260429210808.png}{Pasted image 20260429210808}
 Notemos que en este caso. Los abiertos de la topologia relativa, si estan en la topologia dada por el homeo, pero ademas hay otros, como el ejemplo que di recien, a eso se refiere con que puede ser mas fina.

\end{bookremark}
\phantomsection\label{blk:a7e6ec}
\section{Unicidad de subvariedad dada topologia}
\begin{bookproposition}{Unicidad de subvariedad via inclusion}
Sea $N$ variedad suave y sea $A\subseteq N$. Supongamos que $A$ tiene una topología $\tau$ no necesariamente la topología relativa. Entonces existe a lo sumo una estructura de variedad diferenciable sobre $A$ compatible con $\tau$ tal que

\[
(A,i)
\]
es subvariedad, donde

\[
i:A\to N
\]
es la inclusión. (De hecho podria no haber ninguna)

\end{bookproposition}
\begin{bookproof}{}
\begin{enumerate}
\item Sean $\mathcal F_1,\mathcal F_2$ dos estructuras de variedad diferenciable sobre $A$ compatibles con la topología $\tau$.
\item Para ver que son iguales, alcanza con ver que
\end{enumerate}
\[
\operatorname{Id}:(A,\mathcal F_1)\to (A,\mathcal F_2)
\]
\[
\operatorname{Id}:(A,\mathcal F_2)\to (A,\mathcal F_1)
\]
es difeomorfismo. 3. Usemos el lema de factorizacion 
\bookimage{assets/pasted-image-20260429215527.png}{Pasted image 20260429215527}
 4. Como $F=i$, es claro que la (unica) $\widehat{F}$ que cumple \hyperref[blk:1258d0]{Lema de factorizacion} es

\[
\widehat F=\operatorname{Id}:(A,\mathcal F_1)\to(A,\mathcal F_2)
\]
\begin{enumerate}[start=5]
\item Como la topología que tienen $(A,\mathcal F_1)$ y $(A,\mathcal F_2)$ es la misma, tenemos que $\widehat F=\operatorname{Id}$ es continua, y por tanto $\widehat F$ es suave.
\item Cambiando los roles, también tenemos que $\operatorname{Id}:(A,\mathcal F_2)\to(A,\mathcal F_1)$ es suave.
\item Por tanto, $\operatorname{Id}:(A,\mathcal F_1)\to(A,\mathcal F_2)$ es un difeomorfismo, y entonces $\mathcal F_1=\mathcal F_2$
\item Se necesita que sea difeomorfismo para comparar cartas de ambas estructuras: si $(U,\varphi)\in\mathcal F_1$ y $(V,\psi)\in\mathcal F_2$, entonces en la intersección el cambio de coordenadas se puede escribir usando la identidad como
\end{enumerate}
\[
\psi\circ\varphi^{-1}=\psi\circ\operatorname{Id}\circ\varphi^{-1}.
\]
Como $\operatorname{Id}:(A,\mathcal F_1)\to(A,\mathcal F_2)$ es suave, esta composición es suave. Cambiando los roles se obtiene la suavidad de $\varphi\circ\psi^{-1}$, y por eso las cartas son compatibles.

\bookqed
\end{bookproof}
\phantomsection\label{blk:1933b9}
\begin{bookremark}{}
Notar que por lo que dijimos en \hyperref[blk:a7e6ec]{Teórico 11} entonces para que $\iota: (A,\tau)\rightarrow (N,\tilde{\tau})$ sea continua automaticamente necesitamos que $\tau$ sea un refinamiento de la $\tilde{\tau}$ relativa a $A$. Osea que cualquier abierto de $\tilde{\tau}$ intersecado con $A$ sea abierto de $A$. Al fin y al cabo $\iota ^{-1}(O)=O\cap A$ para cualquier $O \in  \tilde{\tau}$

\end{bookremark}
\begin{bookexample}{Ejemplo donde no hay ninguna}
Si consideramos la variedad ocho osea $(O,f^{-1}|_{(0,2\pi)})$ con $f(t)=(\frac{\sin 2t}{t},\sin t)$ con su topologia $\tau_{1}$  dada por el homemorfismo $f$. Y luego consideramos el mismo conjunto $O$ pero con la topologia $\tau_{2}$ dada por el homeomorfismo $f|_{(-\pi,\pi)}$ Luego

\[
\iota : (O,\tau_{1})\rightarrow (O,\tau_{2})
\]
nunca va a ser continua por que hay abiertos en $\tau_{2}$ que no lo son en $\tau_{1}$ Y esto es independientemente de que estructura le demos a $(O,\tau_{2})$. Por que el punto es que como ya $\iota$ no podra ser continua, tampoco va a ser suave

\end{bookexample}
\begin{bookexample}{Parece pero no es contra ejemplo}
Recordemos que las cartas $(\text{ocho},f^{-1}|_{(0,2\pi)})$ y $(\text{ocho},f^{-1}|_{(-\pi,\pi)})$ nos dan estructuras diferenciables del ocho y luego $(\text{ocho},\iota)$ las hacen subvariedad de $\mathbb{R}^{2}$ Pero esto no contradice \hyperref[blk:1933b9]{Teórico 11} por que estos dos estructuras que le damos a los ochos tienen distintas topologias, que son las heredadas del homeomorfismo $f$. 
\bookimage{assets/pasted-image-20260429220320.png}{Pasted image 20260429220320}
 son abiertos de cada topologia que no esta en la otra

\end{bookexample}
\section{Unicidad de estructura tal que es incrustada}
\begin{bookcorollary}{}
Si $A$ tiene la topología relativa de $N$ y si esta estructura admite una estructura de variedad diferenciable tal que $(A,i)$ es subvariedad de $N$, entonces tal estructura es única y $(A,i)$ es una incrustación.

\end{bookcorollary}
\begin{bookexample}{}
Vimos que

\[
i:S^1\to\mathbb R^2
\]
es una incrustación.

Entonces la estructura de variedad diferencial de $S^2$ es la única estructura de variedad diferenciable sobre $S^{2}$ compatible con la topología heredada en $S^{2}$ de $\mathbb R^3$ que la vuelve subvariedad

\end{bookexample}
